\subsection{The fine-scale Fourier and renewal branch}

This subsection reconstructs the branch left open after
Proposition~\ref{prop:Tao-v7-Prop52-cn-repair}.  The controlling text is the
version-pinned arXiv v7 source, lines~926--1860
\cite{Tao2026v7}.  ArXiv v5 and the published paper are comparison
manifestations, not silent replacements for v7.  Exact source hashes,
line-level defects, and the proof calculations abbreviated here are recorded
in the F024 Fourier--renewal audit and the F028 whole-version audit in the
external QA register.

\begin{sourceissue}[version boundary in the many-triangles estimate]
\label{sourceissue:Tao-v5-v7-many-triangles}
ArXiv v5, source lines~1654--1659, and the journal's Lemma~7.9 state the
unconditional estimate
\begin{equation}
\label{eq:Tao-v5-rip}
 \mathbb E\exp\left(
 -\sum_{p=1}^{t_{\min(r,R)}}1_W(X_p)
 +\varepsilon\min(r,R)\right)\le e^\varepsilon.
\end{equation}
ArXiv v7, source lines~1660--1665, states only
\begin{equation}
\label{eq:Tao-v7-rip}
 \mathbb E1_{\{r\ge R\}}
 \exp\left(-\sum_{p=1}^{t_R}1_W(X_p)+\varepsilon R\right)
 \le e^\varepsilon.
\end{equation}
On \(r\ge R\) the integrands agree, but the v7 formula deletes every path
with \(r<R\); the two statements are not identical.  Moreover, the displayed
v5/journal induction crosses the first triangle's top after \(t_1\).  On the
event \(r=1<R\), the left side of \eqref{eq:Tao-v5-rip} stops at \(t_1\), so
the additional factors used by that induction are absent.  We therefore
record the stronger statement as printed but do not use it.  The weaker v7
statement is proved below.
\end{sourceissue}

\begin{sourceissue}[Section 6 reserve is a substantive version repair]
\label{sourceissue:Tao-v5-v7-section-six-reserve}
Write \(a=\log(4/3)\), \(b=\log2\), and
\[
 q=\frac{b^2}{4a}=0.4175208154\ldots.
\]
The optimized square-root fluctuation cost in Corollary~6.3 is
\(qC_A^2\log N\).  ArXiv v5 and journal equations~(6.7)--(6.8)
retain only \(\frac12bC_A^2\log N=0.3465735903\ldots C_A^2\log N\);
the v5 calculation also drops the factor \(b\) when converting \(2^l\)
to an exponential.  Its printed size estimate therefore has the wrong
leading sign.  An explicit family satisfying the broad v5 concentration
conditions makes one of its natural-number summands exceed \(3^N\); the
construction and its exact scope are recorded in F028.

V7 uses \(0.99bC_A^2\log N=0.6862157088\ldots C_A^2\log N\), and the
stopping event supplies precisely that tightened window after enlarging
\(C_A\).  More generally, a retained coefficient \(\rho\) closes this
calculation exactly when
\[
 \rho>\frac{\log2}{4\log(4/3)}
 =0.6023552099\ldots.
\]
The v5 failure is a defect in the printed proof.  It neither produces two
colliding coordinate tuples nor proves the separation corollary false by a
different argument.
\end{sourceissue}

\begin{proposition}[exact Section 6 reserve boundary]
\label{prop:Tao-Section6-reserve-repair}
Put
\[
 a=\log(4/3),\qquad b=\log2,\qquad
 \rho_*=\frac{b}{4a}.
\]
Let \(n\ge2\), \(L=\log n\), \(C_A>0\), and let
\(a_1,\ldots,a_K\) be positive integers satisfying
\begin{equation}
\label{eq:Tao-Section6-prefix-lower}
 a_{[1,j]}\ge2j-C_A(\sqrt{jL}+L)
 \qquad(1\le j\le K).
\end{equation}
If
\[
 l=a_{[1,K]}
 \le\frac{n\log3}{b}-\rho C_A^2L
\]
for a fixed \(\rho>\rho_*\), then, after first choosing \(C_A\)
sufficiently large in terms of \(\rho-\rho_*\) and then \(n\) sufficiently
large in terms of \(C_A,\rho\),
\begin{equation}
\label{eq:Tao-Section6-Archimedean-separation}
 \sum_{j=1}^{K}3^{j-1}2^{\,l-a_{[1,j]}}<3^n.
\end{equation}
Here
\[
 \rho_*=0.6023552099\ldots,\qquad 0.99>\rho_*.
\]

For the v5 value \(\rho=\frac12\), the universal conclusion
\eqref{eq:Tao-Section6-Archimedean-separation} is false even on sequences
satisfying the source's all-interval concentration inequalities.  By
contrast, on the v7 stopping event \(E_k\cap B_k\), \(C_A\ge600\) gives
\[
 a_{[1,k+1]}
 \le\frac{n\log3}{\log2}-0.99C_A^2\log n.
\]
\end{proposition}

\begin{proof}
From \eqref{eq:Tao-Section6-prefix-lower},
\[
\begin{split}
 S
 :=\sum_{j=1}^{K}3^{j-1}2^{\,l-a_{[1,j]}}
 &\ll
 3^n e^{-\rho bC_A^2L+bC_AL}
 \sum_{j\ge1}e^{-aj+bC_A\sqrt{jL}}.
\end{split}
\]
The exact completion of the square is
\[
 -aj+bC_A\sqrt{jL}
 =-a\left(\sqrt j-\frac{bC_A\sqrt L}{2a}\right)^2
 +\frac{b^2}{4a}C_A^2L.
\]
Splitting the sum into unit intervals around its maximum, or comparing it
with the corresponding Gaussian integral after \(u=\sqrt j\), gives
\[
 \sum_{j\ge1}e^{-aj+bC_A\sqrt{jL}}
 \ll(1+C_A^2L)e^{b^2C_A^2L/(4a)}.
\]
The coefficient of \(C_A^2L\) in the resulting exponent is
\[
 -\rho b+\frac{b^2}{4a}=-b(\rho-\rho_*).
\]
It is strictly negative under the stated hypothesis.  It absorbs first the
linear \(bC_AL\) term after \(C_A\) is chosen, and then the polynomial factor
after \(n\) is chosen.  This proves \(S<3^n\).

For the v5 counterexample, set
\[
 \lambda=\frac{2a}{b},\qquad \gamma=\lambda^{-2},\qquad
 J=\lfloor\gamma C_A^2L\rfloor,
\]
and, for \(1\le i\le J\), put
\[
 d_i=\lfloor\lambda i\rfloor-\lfloor\lambda(i-1)\rfloor,\qquad
 a_i=2-d_i;
\]
put \(d_i=0\) and \(a_i=2\) afterwards.  Thus \(a_i\in\{1,2\}\), and the
deficit on every interval of length \(r\) is at most \(\lambda r+1\).
For \(r\le J\), \(\lambda r\le C_A\sqrt{rL}\); for \(r>J\), the entire
deficit is at most \(C_A\sqrt{rL}+O(1)\).  Hence the source's additional
\(C_AL\) tolerance makes every printed interval inequality valid.

Let \(D=\lfloor\lambda J\rfloor\).  Choose \(l\), with \(l+D\) even, within
two of
\[
 \frac{n\log3}{b}-\frac12C_A^2L,
\]
and put \(K=(l+D)/2\).  For sufficiently large \(n\), \(K\ge J\) and
\(a_{[1,K]}=2K-D=l\).  The \(j=J\) summand satisfies
\[
 \log\frac{3^{J-1}2^{\,l-a_{[1,J]}}}{3^n}
 =
 \left(\frac{b^2}{4a}-\frac12b\right)C_A^2L+O(1).
\]
The coefficient is
\[
 0.4175208154\ldots-0.3465735903\ldots
 =0.0709472252\ldots>0,
\]
so that single summand eventually exceeds \(3^n\).

Finally put \(T=n\log3/b-C_A^2L\).  On \(B_k\),
\(a_{[1,k]}\le T<a_{[1,k+1]}\).  The \(i=k,j=k+1\) inequality on \(E_k\)
and \(a_k\ge1\) yield
\[
 a_{[1,k+1]}\le T+1+C_A(\sqrt L+L).
\]
For \(n\ge2\), the last overshoot is at most \(6C_AL\), and for
\(C_A\ge600\) this is at most \(0.01C_A^2L\).  This gives the v7
\(0.99\)-window.
\end{proof}

\begin{nonclaim}
Proposition~\ref{prop:Tao-Section6-reserve-repair} refutes the universal
Archimedean size assertion in the printed v5 proof, not the separation
corollary by every possible proof.  The construction supplies no pair of
colliding tuples.  The coefficient \(0.99\) is convenient rather than
canonical, and the threshold \(\rho_*\) belongs to this exact
optimization-and-absorption argument.  No Fourier-decay theorem, main
theorem, or Collatz endpoint follows from this finite margin calculation.
\end{nonclaim}

\begin{definition}[uniform lifting and fibre averaging]
\label{def:Tao-uniform-lift}
For \(0\le r\le N\), let
\[
 L_{r\to N}:\mathbb R^{\mathbb Z/3^r\mathbb Z}
 \longrightarrow\mathbb R^{\mathbb Z/3^N\mathbb Z},
 \qquad
 (L_{r\to N}\mu)(y)=3^{r-N}\mu(y\bmod3^r).
\]
For the reduction map
\(\pi_{N,r}:\mathbb Z/3^N\mathbb Z\to\mathbb Z/3^r\mathbb Z\), its fibre
above \(x\) is \(x+3^r\mathbb Z/3^N\mathbb Z\).  On signed measures on the
larger group define the fibre-constant projection
\[
 P_{r,N}:\mathbb R^{\mathbb Z/3^N\mathbb Z}
 \longrightarrow\mathbb R^{\mathbb Z/3^N\mathbb Z},\qquad
 (P_{r,N}\nu)(y)=3^{r-N}
 \sum_{\substack{z\bmod3^N\\z\equiv y\ (3^r)}}\nu(z).
\]
\end{definition}

Both maps preserve total mass, \(P_{r,N}\) is an \(\ell^1\)-contraction,
and \(L_{r\to N}\) is an \(\ell^1\)-isometry.  If \(\mu_N\) is the law of
Tao's \(N\)-Syracuse random variable, projective consistency says
\((\pi_{N,r})_*\mu_N=\mu_r\).  Substitution in the two displayed maps gives
the typed identity
\begin{equation}
\label{eq:Tao-lift-projective}
 P_{r,N}\mu_N=L_{r\to N}\mu_r.
\end{equation}

\begin{proposition}[exact reduction from characteristic decay to fine-scale mixing]
\label{prop:Tao-v7-Fourier-to-mixing}
For every \(q\in\Npos\), let \(\mu_q\) be the law on
\(\mathbb Z/3^q\mathbb Z\) just defined.  Assume the v7
characteristic-function estimate, uniformly for
\(\xi\in\mathbb Z/3^q\mathbb Z\) with \(3\nmid\xi\),
\begin{equation}
\label{eq:Tao-v7-characteristic-input}
 \left|\mathbb E
 e^{-2\pi i\xi\operatorname{Syrac}(\mathbb Z/3^q\mathbb Z)/3^q}
 \right|\ll_Bq^{-B}
\end{equation}
for every fixed \(B>0\), with an implied constant depending only on \(B\).
After correcting source line~1051 from \(3^q\)
to \(3^{q-1}\), and line~1093 from \(m\) to \(k+1\), the source argument
proves
\[
 \|\mu_N-L_{r\to N}\mu_r\|_1\ll_A r^{-A}
 \qquad(N,r\in\Npos,\ 1\le r\le N).
\]
This is exactly Proposition~1.14, including its domain and uniformity.
\end{proposition}

\begin{proof}
First suppose \(0.9N\le r\le N\).  Let \(E\) be the source's all-interval
concentration event.  For a fixed \(k\), let \(E_k\supseteq E\) retain only
the concentration inequalities involving \(a_1,\ldots,a_{k+1}\).  The source's
concentration estimate gives
\(\mathbb P(E_k\setminus E)=O_{A,C_A}(N^{-A-1})\), so the triangle
inequality permits \(E\) to be replaced by \(E_k\) at the stated error
scale.  On \(E_k\), stop at the unique \(k\) satisfying
\[
 a_{[1,k]}\le \frac{N\log3}{\log2}-C_A^2\log N
 <a_{[1,k+1]},
\]
and condition also on \(a_{[1,k+1]}=l\).  Write \(B_k\) for this stopping
event and \(C_{k,l}=\{a_{[1,k+1]}=l\}\).  Put
\[
 T_N=\frac{N\log3}{\log2}-C_A^2\log N.
\]
The \(i=k,j=k+1\) inequality in \(E_k\), together with \(a_k\ge1\), gives
\[
 a_{k+1}\le1+C_A(\sqrt{\log N}+\log N),
\]
and hence
\begin{equation}
\label{eq:Tao-v7-point99-window}
 l=a_{[1,k+1]}
 \le T_N+1+C_A(\sqrt{\log N}+\log N)
 \le\frac{N\log3}{\log2}-0.99C_A^2\log N
\end{equation}
after choosing \(C_A\ge600\) and then \(N\ge2\).  The last displayed
sufficient choice follows from
\(\sqrt{\log N}\le2\log N\), \(2\le3\log N\), and
\(6C_A\le0.01C_A^2\).  The prefix and suffix offset maps then have the
exact typed decomposition
\[
 X_N=F_{k+1}(a_{k+1},\ldots,a_1)
 +3^{k+1}2^{-l}F_{N-k-1}(a_N,\ldots,a_{k+2})
 \pmod{3^N}.
\]
The events \(E_k,B_k,C_{k,l}\) and the prefix offset depend only on
\(a_1,\ldots,a_{k+1}\), and hence are independent of the suffix variables.
Let
\(\eta\in\mathbb Z/3^N\mathbb Z\) be a frequency surviving fibre averaging,
 so \(\eta\notin3^{N-r}\mathbb Z/3^N\mathbb Z\).  There are unique
\[
 0\le v<N-r,
 \qquad \bar\eta'\in(\mathbb Z/3^{N-v}\mathbb Z)^\times,
\]
where \(v=\nu_3(\eta)\) and
\[
 \eta/3^v=2^l\bar\eta'\pmod{3^{N-v}}.
\]
Thus the quotient is asserted unique only in the correct quotient ring, not
as a residue modulo \(3^N\).  Put
\[
 q=N-k-v-1.
\]
The concentration range for \(k\), together with \(v<N-r\le0.1N\), gives
\[
 q\ge0.9N-\frac{N\log3}{2\log2}
       -O(C_A\sqrt{N\log N})-1\gg N
\]
for sufficiently large \(N\).  Let
\(\eta'_q\in(\mathbb Z/3^q\mathbb Z)^\times\) be the reduction of
\(\bar\eta'\).  Cancelling the exact powers of three and two in the suffix
character gives the typed frequency morphism
\begin{align*}
 &\mathbb E\exp\left(
 -\frac{2\pi i\eta\,3^{k+1}2^{-l}
 F_{N-k-1}(a_N,\ldots,a_{k+2})}{3^N}\right)\\
 &\qquad=
 \mathbb E\exp\left(-\frac{2\pi i\eta'_q
 F_{N-k-1}(a_N,\ldots,a_{k+2})}{3^q}\right)\\
 &\qquad=
 \mathbb E\exp\left(-\frac{2\pi i\eta'_q
 \operatorname{Syrac}(\mathbb Z/3^q\mathbb Z)}{3^q}\right).
\end{align*}
The final equality is projective consistency of the offset law from
\(N-k-1\) to \(q\), and
\eqref{eq:Tao-v7-characteristic-input} bounds this factor by
\(O_B(N^{-B})\), uniformly in every surviving frequency and admissible
\((k,l)\).

Write \(\mathcal E_{r,N;k,l}^2\) for the squared oscillation at this fixed
pair \((k,l)\).  Fourier inversion, removal of the frequencies divisible by
\(3^{N-r}\), the preceding suffix bound, and Plancherel give
\begin{equation}
\label{eq:Tao-prefix-collision}
 \mathcal E_{r,N;k,l}^2\ll_B N^{-2B}3^N\sum_y p_y^2,
\end{equation}
where, with the source events \(E_k,B_k,C_{k,l}\),
\[
 p_y=\mathbb P\bigl(
 F_{k+1}(a_{k+1},\ldots,a_1)\equiv y\pmod{3^N},
 E_k\cap B_k\cap C_{k,l}\bigr).
\]
Thus \((p_y)_y\) is the unnormalised restricted joint mass, not a
conditional probability distribution.

For positive coordinates,
\[
 F_q(a_1,\ldots,a_q)
 =\sum_{u=1}^q3^{q-u}2^{-a_{[u,q]}}.
\]
The unique term of least 2-adic valuation is
\(3^{q-1}2^{-a_{[1,q]}}\).  Hence equality of two rational offsets forces
equality of their total valuation sums.  Subtraction gives equality of the
two \((q-1)\)-coordinate tails, and induction proves that \(F_q\) has
singleton fibres on \(\Npos^q\).

Under the fixed-total and concentration conditions, multiplying a congruence
of two prefix offsets by \(2^l\) gives a congruence modulo \(3^N\) between
two positive integers.  Their prefix inequalities and
\eqref{eq:Tao-v7-point99-window} meet the hypotheses of
Proposition~\ref{prop:Tao-Section6-reserve-repair} with \(\rho=0.99\).
Hence both positive integers are strictly below \(3^N\).  They are therefore
equal as integers.  Dividing by \(2^l\) and applying rational injectivity
proves residue-level separation.
Thus
\[
 0\le p_y\le2^{-l},\qquad \sum_y p_y\le1,
\]
and consequently
\[
 3^N\sum_yp_y^2\le3^N(\max_yp_y)\sum_yp_y
 \le3^N2^{-l}=N^{O(C_A^2)}.
\]
Taking \(B\) sufficiently large in
\eqref{eq:Tao-v7-characteristic-input} absorbs this precise polynomial and
closes \eqref{eq:Tao-prefix-collision} with an arbitrary prescribed
\(N^{-2A-4}\) reserve.  There are \(O_A(N)\) admissible stopping indices
\(k\), and, from the displayed fixed-total window, \(O_A(\log N)\)
admissible totals \(l\).  Cauchy--Schwarz converts the squared bound back to
the fixed-pair \(\ell^1\)-oscillation, and the union bound over these
\((k,l)\) uses less than the two reserved powers of \(N\).  The complement
of the all-interval concentration event is \(O_A(N^{-A-1})\) after enlarging
\(C_A\).  Adding these disjoint fixed-pair contributions proves the
restricted \(\ell^1\) estimate, rather than silently identifying a
conditioned estimate with the unconditional one.

For general \(10\le r\le N\), construct
\(N=n_0>n_1>\cdots>n_t=r\) by setting
\(n_{i+1}=\max(r,\lceil0.9n_i\rceil)\).  Because every \(n_i\ge10\)
before termination, this sequence decreases strictly, and every adjacent
pair is in the proved range.  By \eqref{eq:Tao-lift-projective}, functoriality of the
lifts, and preservation of the \(\ell^1\)-norm,
\[
\begin{aligned}
 \|\mu_N-L_{r\to N}\mu_r\|_1
 &\le\sum_{i=0}^{t-1}
 \|L_{n_i\to N}\mu_{n_i}
       -L_{n_{i+1}\to N}\mu_{n_{i+1}}\|_1\\
 &=\sum_{i=0}^{t-1}
 \|\mu_{n_i}-L_{n_{i+1}\to n_i}\mu_{n_{i+1}}\|_1.
\end{aligned}
\]
Applying the restricted estimate with exponent \(A+1\), and summing the
geometrically separated levels, gives \(O_A(r^{-A})\).  This supplies the
lifting sequence implicit in source lines~932--936.  When \(1\le r<10\),
the triangle inequality gives a bound of two, which is
\(O_A(r^{-A})\) after enlarging the constant by at most \(2\,9^A\).
The finitely many \(N\) excluded while proving the restricted large-\(N\)
estimate are absorbed in the same \(A\)-dependent constant.  These two
finite-scale clauses are separate from the strictly decreasing chain.
\end{proof}

For \(d\in\Npos\), \(t\in\mathbb R_{\ge0}\), and \(y\in\mathbb R^d\),
we use Tao's Gaussian--exponential kernel
\begin{equation}
\label{eq:Tao-G-kernel}
 G_t(y)=
 \begin{cases}
  \exp(-\|y\|^2/t)+\exp(-\|y\|),&t>0,\\
  \exp(-\|y\|),&t=0.
 \end{cases}
\end{equation}

\begin{lemma}[the corrected two-dimensional local factor]
\label{lem:Tao-v7-local-factor}
In the proof of Tao's Lemma~2.2, the source-line-516 factor is
\(n^{-d/2}\), not \(n^{-1/2}\).  In particular, let
\(\mathbf H_1,\mathbf H_2,\ldots\) be iid copies of the holding vector
\(\mathbf H\in\mathbb Z^2\) of
Proposition~\ref{prop:Tao-v7-pairing-holding}, put
\(\mathbf H_{[1,q]}=\sum_{i=1}^q\mathbf H_i\) with the empty sum at
\(q=0\), and recall \(\mathbb E\mathbf H=(4,16)\).  There exist constants
\(c,C>0\), depending only on the law of \(\mathbf H\), such that for every
\(q\in\Nzero\) and \((x,t)\in\mathbb Z^2\),
\[
 \mathbb P(\mathbf H_{[1,q]}=(x,t))
 \le C(q+1)^{-1}G_q(c((x,t)-q(4,16))).
\]
\end{lemma}

\begin{proof}
After the source's contour shift, the Fourier integral contains
\(e^{-cn|u|^2}\) on a bounded subset of \(\mathbb R^d\).  The change of
variables \(v=\sqrt n\,u\) gives
\[
 \int e^{-cn|u|^2}\,du
 \le n^{-d/2}\int_{\mathbb R^d}e^{-c|v|^2}\,dv.
\]
All remaining source steps preserve this dimension.  Setting \(d=2\)
gives the displayed factor.
\end{proof}

\begin{proposition}[pairing, phase, and exact holding vector]
\label{prop:Tao-v7-pairing-holding}
Let \(n\in\Npos\), let
\(\xi\in\mathbb Z/3^n\mathbb Z\) satisfy \(3\nmid\xi\), and define
\[
 \chi:\mathbb Z[1/2]\longrightarrow\mathbb C,\qquad
 \chi(x)=\exp\left(-\frac{2\pi i\xi(x\bmod3^n)}{3^n}\right).
\]
Fix \(0<\varepsilon\le e^{-10}\).  For
\((r,l)\in\Npos\times\mathbb Z\), define the
signed fractional phase
\[
 \theta(r,l)=\left\{
 \frac{\xi3^{2r-2}(2^{-l+1}\bmod3^n)}{3^n}
 \right\}\in(-1/2,1/2],
\]
and define the typed partition of the strip
\[
 B=\{(r,l):1\le r\le\lfloor n/2\rfloor,
                  |\theta(r,l)|\le\varepsilon\},\qquad
 W=\bigl(\{1,\ldots,\lfloor n/2\rfloor\}\times\mathbb Z\bigr)\setminus B.
\]
Let \((a_i)_{i\ge1}\) be iid with
\(\mathbb P(a_i=u)=2^{-u}\) for \(u\in\Npos\), and put
\[
 S_\chi(m)=\mathbb E\chi\left(
 \sum_{v=1}^{m}3^{v-1}2^{-\sum_{w=1}^{v}a_w}
 \right)\qquad(m\in\Nzero),
\]
with the empty-sum convention \(S_\chi(0)=1\).
For independent copies \(a_1,a_2\) of this geometric variable, define
\[
 \begin{aligned}
 f&:\mathbb Z[1/2]\times\{2,3,\ldots\}\longrightarrow\mathbb C,\\
 f(x,b)&=\mathbb E\!\left[
 \chi\bigl(x(2^{a_2}+3)\bigr)\mid a_1+a_2=b\right],\\
 g&:\mathbb Z[1/2]\longrightarrow\mathbb C,\\
 g(x)&=\mathbb E\chi\bigl(x2^{-a_1}\bigr).
 \end{aligned}
\]
Let \(q=\lfloor n/2\rfloor\),
\(b_r=a_{2r-1}+a_{2r}\), and
\(B_r=\sum_{u=1}^rb_u\), with \(B_0=0\).  Then the \(b_r\) are iid with
\(\mathbb P(b_r=b)=(b-1)2^{-b}\) for \(b\ge2\).  Expectation in the next
two displays is over this joint Pascal law, and conditioning on the
\(b_r\) gives the exact formulas
\[
 S_\chi(2q)=\mathbb E
 \prod_{r=1}^qf(3^{2r-2}2^{-B_r},b_r)
\]
and
\[
 S_\chi(2q+1)=\mathbb E\left[
 \prod_{r=1}^qf(3^{2r-2}2^{-B_r},b_r)\,
 g(3^{2q}2^{-B_q})\right].
\]
At every white point \((r,l)\),
\[
 |f(3^{2r-2}2^{-l},3)|\le e^{-\varepsilon^3}.
\]
The occurrences \(b_r=3\) form a renewal process with iid holding vector
\[
 \mathbf H=(1,3)+\sum_{i=1}^{R}(1,Q_i),
 \qquad
 \mathbb P(R=s)=\frac14\left(\frac34\right)^s
 \qquad(s\in\Nzero),
\]
where \(R\) is independent of the iid \(Q_i\), and the \(Q_i\) have the
Pascal law conditioned on \(Q_i\ne3\).  This vector
has exponential tail, full lattice span, and mean \((4,16)\).
\end{proposition}

\begin{proof}
The pairing identities follow by conditioning while retaining the random
odd-\(n\) terminal factor inside the expectation.  Conditional on \(b=3\),
the fibres \((a_1,a_2)=(1,2),(2,1)\) have equal mass.  Therefore
\[
 f(x,3)=\frac{\chi(5x)+\chi(7x)}2
 =\frac{\chi(5x)}2(1+\chi(2x)).
\]
For \(x=3^{2r-2}2^{-l}\), the corrected character identity is
\(\chi(2x)=e^{-2\pi i\theta(r,l)}\).  Since
\(\theta(r,l)\in(-1/2,1/2]\),
\[
 |f(x,3)|=\cos(\pi\theta(r,l)).
\]
The elementary bound
\(\cos(\pi u)\le e^{-2u^2}\) for \(|u|\le1/2\) gives, at a white point,
\[
 |f(x,3)|\le e^{-2\varepsilon^2}\le e^{-\varepsilon^3},
\]
where the last inequality uses \(0<\varepsilon\le e^{-10}<2\).

The success probability is \(\mathbb P(b=3)=1/4\), giving the displayed
failure-before-success representation.  For
\(k=(k_1,k_2)\in\mathbb R^2\), put
\(M_Q(k)=\mathbb E e^{k_1+k_2Q}\).  On the open neighbourhood of zero on
which \(M_Q(k)<\infty\) and \(3M_Q(k)<4\), the normalized moment generating
function is
\begin{equation}
\label{eq:Tao-hold-mgf-corrected}
 \mathbb E e^{\mathbf H\cdot k}
 =\sum_{s=0}^\infty\frac14\left(\frac34\right)^s
 e^{(1,3)\cdot k}M_Q(k)^s
 =\frac{e^{(1,3)\cdot k}}{4-3M_Q(k)}.
\end{equation}
It equals one at \(k=0\) and is finite near zero.  This corrects the
off-by-one series at source line~1408 and proves the exponential tail.
The support contains \((1,3),(2,5),(2,7),(2,8)\); its differences generate
\((0,1)\) and \((1,0)\), so it is contained in no coset of a proper subgroup
of \(\mathbb Z^2\).  Finally,
\(\mathbb E Q=13/3\), \(\mathbb ER=3\), and direct substitution gives
\(\mathbb E\mathbf H=(4,16)\).
\end{proof}

\begin{proposition}[corrected black-triangle decomposition]
\label{prop:Tao-v7-black-triangles}
Retain \(n,\xi,\varepsilon,\theta,B,W\) from
Proposition~\ref{prop:Tao-v7-pairing-holding}.
After the forced corrections at v7 lines~1272, 1285--1287, and
1311--1313, \(B\) is a disjoint union of lattice triangles
\[
 \Delta=\{(j,l):j\ge j_\Delta,\ l\le l_\Delta,\
 (j-j_\Delta)\log9+(l_\Delta-l)\log2\le s_\Delta\}.
\]
Here \(j_\Delta\in\Npos\), \(l_\Delta\in\mathbb Z\),
\(s_\Delta\in\mathbb R_{\ge0}\), and the displayed set is intersected with
the lattice \(\Npos\times\mathbb Z\).
With \(\delta=\frac1{10}\log(1/\varepsilon)\), every triangle satisfies
\[
 \Delta\subseteq
 \{(j,l)\in\Npos\times\mathbb Z:j\le n/2-\delta\},
\]
and
distinct triangles have Euclidean distance at least \(\delta\).
\end{proposition}

\begin{proof}
The exact recurrences are
\[
 \theta(j+1,l)\equiv9\theta(j,l),\qquad
 \theta(j,l-1)\equiv2\theta(j,l)\pmod{\mathbb Z}.
\]
Call a point weakly black when \(|\theta|\le1/100\).  The recurrences give
three exact propagation rules.  First, if \((j,l)\) is weakly black and
\((j+1,l)\) is black, then \(|9\theta(j,l)|\le9/100<1/2\), so there is no
wrap and \(|\theta(j,l)|=|\theta(j+1,l)|/9\le\varepsilon\).  The same
argument with \(|2\theta(j,l)|\le2/100<1/2\) applies when \((j,l-1)\) is
black.  Second, if \((j+1,l)\) and \((j,l-1)\) are weakly black, then
\[
 |\theta(j,l)|\le |\theta(j+1,l)|+4|\theta(j,l-1)|\le5/100;
\]
now \(|2\theta(j,l)|\le1/10<1/2\), so the doubling recurrence has no wrap
and \(|\theta(j,l)|=|\theta(j,l-1)|/2\le1/200\).  Third, if
\((j-1,l)\) and \((j,l-1)\) are weakly black, the factor-nine recurrence
first gives \(|\theta(j,l)|\le9/100\); hence
\(|2\theta(j,l)|\le18/100<1/2\), and doubling again gives
\(|\theta(j,l)|\le1/200\).

If \((j,l)\) is black, multiplying its phase to the primitive residue at
level \(3^n\) gives
\[
 3^{n+1-2j}\varepsilon\ge\frac13.
\]
Taking logarithms gives
\[
 j\le \frac n2+1-\frac{\log(1/\varepsilon)}{2\log3}
 \le \frac n2-\frac1{10}\log(1/\varepsilon)
\]
for \(\varepsilon\le e^{-10}\).  There is no infinite upward black run:
if \((j,l')\) were black for every \(l'\ge l\), repeated nonwrapping
doubling would give
\(|\theta(j,l')|\le2^{l-l'}\varepsilon\), contradicting the same primitive
residue lower bound as \(l'\to\infty\).

Starting from a black \((j,l)\), let \(l_*\ge l\) be the unique top of its
consecutive black vertical run, and let \(j_*\le j\) be the unique left end
of the consecutive black run on row \(l_*\).  Existence and uniqueness follow
from the preceding no-infinite-run statement and the lower boundary
\(j=1\).  Because \(2j_*-2<n\), while neither \(\xi\) nor \(2\) is
divisible by three, the defining rational phase is not integral; hence
\(\theta(j_*,l_*)\ne0\).  Write
\(|\theta(j_*,l_*)|=\varepsilon e^{-s_*}\).  On the corrected domain
\(j'\ge j_*\), \(l'\le l_*\), multiplication by
\(9^{j'-j_*}2^{l_*-l'}\) gives
\[
 |\theta(j',l')|\le
 \varepsilon e^{-s_*+(j'-j_*)\log9+(l_*-l')\log2},
\]
with equality before the right side reaches \(1/2\).  Thus all points of the
anchored triangle are black.  If an exterior Case~1 point is within
\(\delta\), and
\(X=-s_*+(j'-j_*)\log9+(l_*-l')\log2\), then
\[
 \varepsilon<\varepsilon e^X
 \le\varepsilon^{1-(\log9+\log2)/10}<\tfrac12,
\]
so \(|\theta(j',l')|=\varepsilon e^X>\varepsilon\).  The corrected Case~2
argument is as follows.  If \(l'>l_*\) and a nearby exterior point
\((j',l')\) were black, Euclidean distance at most \(\delta\) gives
\[
 0<(l'-l_*)\log2\le\delta\log2,
 \qquad
 (j'-j_*)\log9\le s_*+\delta\log9.
\]
Exact downward phase transport then gives
\[
 |\theta(j',l_*+1)|
 \le\varepsilon^{1-\log2/10}<1/100.
\]
When \(j'>j\), the anchored phase bound also gives
\[
 |\theta(j'-1,l_*)|
 \le\varepsilon^{1-\log9/10}<1/100.
\]
The second propagation rule iterates left from \(j'-1\) to \(j\).  If
\(j'=j\), \((j,l_*+1)\) is already the weakly black point supplied by the
first display.  Thus in either case \((j,l_*+1)\) is weakly black, and
blackness of \((j,l_*)\) makes it black by the first rule, contradicting the
definition of \(l_*\).  When \(j'<j\), the black row from \(j_*\) through
\(j\) and the third rule instead propagate weak blackness right from
\((j',l_*+1)\) to \((j,l_*+1)\), giving the same contradiction.

In Case~3, \(j'<j_*\).  The same distance condition gives
\[
 |l'-l_*|\log2\le s_*+\delta\log2,
 \qquad 0<(j_*-j')\log9\le\delta\log9.
\]
Assuming \((j',l')\) black, exact horizontal phase transport gives
\[
 |\theta(j_*-1,l')|
 \le\varepsilon^{1-\log9/10}<1/100.
\]
If \(l'\ge l_*\), vertical transport gives the explicit bound
\[
 |\theta(j_*-1,l_*)|
 \le\varepsilon^{1-(\log9+\log2)/10}<1/100.
\]
If \(l'<l_*\), the anchored bound gives
\(|\theta(j_*,l'+1)|\le\varepsilon^{1-\log2/10}<1/100\), and the second
rule propagates upward until \((j_*-1,l_*)\) is weakly black.  In either
subcase, the first propagation rule and blackness of \((j_*,l_*)\) make
\((j_*-1,l_*)\) black, contradicting the defining minimality of \(j_*\).

Hence every exterior point within \(\delta\) is white.  Since
\(\delta\ge1\), the black horizontal and vertical paths defining the anchor
cannot cross the boundary: the first black path point outside the triangle
would be an exterior black point within distance one.  Thus every original
black point lies in its anchored triangle.

Conversely, every point of this triangle reconstructs the same anchor.
For each column occurring in the triangle, the entire vertical segment up to
row \(l_*\) is black and the point at height \(l_*+1\) is an exterior collar
point, hence white.  On row \(l_*\), the entire horizontal segment back to
\(j_*\) is black and, unless \(j_*=1\), the point at \(j_*-1\) is an
exterior collar point, hence white.  Therefore its reconstructed pair is
exactly \((j_*,l_*)\).  The anchored triangles partition \(B\), and the
white collar makes distinct members \(\delta\)-separated.

Confinement of the entire abstract triangle follows directly from the
anchor denominator, without presupposing confinement of its lattice points.
The nonzero signed phase satisfies
\[
 |\theta(j_*,l_*)|\ge3^{2j_*-2-n}.
\]
Since its other expression is \(\varepsilon e^{-s_*}\),
\[
 s_*\le(n-2j_*+2)\log3-\log(1/\varepsilon).
\]
Consequently its real right tip satisfies
\[
 j_*+\frac{s_*}{\log9}
 \le\frac n2+1-\frac{\log(1/\varepsilon)}{2\log3}
 \le\frac n2-\delta.
\]
Every point of the triangle lies weakly to the left of that tip, proving the
asserted confinement and, in particular, the real-tip bound needed later.
\end{proof}

\begin{lemma}[first-passage Green kernel]
\label{lem:Tao-v7-first-passage-repaired}
Let \((\mathbf H_i)_{i\ge1}\) be iid copies of \(\mathbf H\), let
\(\mathbf V_q=\sum_{i=1}^q\mathbf H_i\), and set
\(\mathbf V_0=(0,0)\).  For \(s\in\Nzero\), put
\(\kappa_s=\inf\{q\in\Npos:(\mathbf V_q)_2>s\}\).  For every
\(j,l\in\mathbb Z\) with \(l>s\),
\[
 \mathbb P(\mathbf V_{\kappa_s}=(j,l))
 \ll\frac{e^{-c(l-s)}}{\sqrt{1+s}}
 G_{1+s}(c(j-s/4)).
\]
Moreover, there is an absolute \(C\ge1\) such that the following typed
post-passage estimate holds.  Let \(m\in\Npos\) be sufficiently large,
let \(s,p\in\Nzero\), let \(r_0\in[1,\infty)\), and let \(h>0\).  Suppose
\[
 s>\frac{m}{\log^2m},\qquad 0\le p\le m^{0.1},\qquad
 4r_0\le h\le m^{0.4},
\]
and let \(\Sigma\subset\mathbb Z\) have pairwise spacing at least \(h\).
Then, for all sufficiently large \(m\),
\[
 \mathbb P\!\left(
 \operatorname{dist}((\mathbf V_{\kappa_s+p})_1,\Sigma)\le r_0
 \right)\le C\frac {r_0}h.
\]
\end{lemma}

\begin{proof}
Set
\(U(x,t)=\sum_{q\ge0}\mathbb P(\mathbf V_q=(x,t))\) for
\((x,t)\in\mathbb Z\times\Nzero\).  By
Lemma~\ref{lem:Tao-v7-local-factor},
\[
 \mathbb P(\mathbf V_q=(x,t))
 \ll(q+1)^{-1}G_q(c((x,t)-q(4,16))).
\]
Writing \(z=x-t/4\), the invertible coordinate identity
\[
 (x-4q,t-16q)=\left(z+\frac{t-16q}{4},t-16q\right)
\]
preserves norms up to two absolute constants.  Put
\(\langle t\rangle=1+t\), and, after harmlessly weakening the denominator
\(q\) in the Gaussian to \(q+1\), set
\[
 A_q(z,t)=\frac1{q+1}
 \exp\left(-c\frac{z^2+(t-16q)^2}{q+1}\right).
\]
In the central range the complete Gaussian calculation is
\[
 \sum_{t/2\le16q\le2t}A_q(z,t)
 \ll \langle t\rangle^{-1/2}
 e^{-c'z^2/\langle t\rangle}.
\]
Indeed \(q+1\asymp\langle t\rangle\) there and the discrete Gaussian in
\(t-16q\) has mass \(O(\sqrt{\langle t\rangle})\).  In the lower range,
\(|t-16q|\gg t\) and \(q+1\ll\langle t\rangle\); in the upper range use
\[
 (q+1)+\frac{z^2}{q+1}
 \gg \langle t\rangle+\min\left(
 \frac{z^2}{\langle t\rangle},|z|\right).
\]
Geometric summation therefore gives, respectively,
\[
 \sum_{16q<t/2}A_q(z,t)
 \ll\langle t\rangle^{-1/2}e^{-c'z^2/\langle t\rangle},
 \qquad
 \sum_{16q>2t}A_q(z,t)
 \ll\langle t\rangle^{-1/2}
 \left(e^{-c'z^2/\langle t\rangle}+e^{-c'|z|}\right).
\]
The exponential half of \(G_q\) is even smaller, since
\[
 \sum_{q\ge0}\frac{e^{-c(|z|+|t-16q|)}}{q+1}
 \ll\langle t\rangle^{-1/2}e^{-c'|z|}.
\]
The cases \(t=0,1\) are absorbed by changing the absolute constants.  Thus
\[
 U(x,t)\ll(1+t)^{-1/2}G_{1+t}(c'(x-t/4)).
\]

At first passage the final vertical increment is at least \(l-s\).  Using
its exponential tail and independence from the earlier increments gives
\[
 \mathbb P(\mathbf V_{\kappa_s}=(j,l))
 \ll e^{-c(l-s)}\sum_{r=0}^s\sum_{u\ge1}
 e^{-c(r+u)}U(j-u,s-r).
\]
For
\[
 K_t(y)=(1+t)^{-1/2}
 \left(e^{-c y^2/(1+t)}+e^{-c|y|}\right),
\]
the two needed convolution inequalities, with absolute
\(a,c_1,c_2>0\), are
\begin{align}
 \sum_{r=0}^{s}e^{-ar}K_{s-r}(x+r/4)
 &\ll K_s(c_1x),\label{eq:Tao-vertical-convolution}\\
 \sum_{u\ge1}e^{-au}K_s(x-u)
 &\ll K_s(c_2x).\label{eq:Tao-horizontal-convolution}
\end{align}
For \eqref{eq:Tao-vertical-convolution}, split at \(r=s/2\).  In the first
part the two variances are comparable and
\[
 ar+c\frac{(x+r/4)^2}{1+s-r}
 \ge c'\frac{x^2}{1+s}+c'r;
\]
in the second part the factor \(e^{-ar}\), together with the retained
\(x+r/4\) decay, obeys the explicit bound
\[
 ar+c\frac{(x+r/4)^2}{1+s-r}
 \ge c'\left(s+\min\left(\frac{x^2}{1+s},|x|\right)\right)
 \qquad(r>s/2).
\]
Since
\(\sum_{r>s/2}(1+s-r)^{-1/2}e^{-c's}\ll
(1+s)^{-1/2}\), this absorbs the lost square-root prefactor.  For the
exponential-kernel half, the corresponding inequalities are
\[
 ar+c|x+r/4|\ge c'(|x|+r)\quad(r\le s/2),
 \qquad
 ar+c|x+r/4|\ge c'(|x|+s)\quad(r>s/2).
\]
For \eqref{eq:Tao-horizontal-convolution}, split at \(u=|x|/2\).  In the
first range \(|x-u|\ge|x|/2\) whenever \(x\ge0\), while for \(x<0\) it
holds for every \(u\ge1\).  In the complementary range the geometric sum
of \(e^{-au}\) is \(O(e^{-a|x|/2})\).  These statements apply separately
to the Gaussian and exponential halves of \(K_s\) and prove the second
displayed inequality.  In the actual summand \(U(j-u,s-r)\), the kernel
coordinate is \(j-s/4-u+r/4\).  Thus first substitute
\(x=j-s/4-u\) in \eqref{eq:Tao-vertical-convolution}; after that summation,
substitute \(x=j-s/4\) in
\eqref{eq:Tao-horizontal-convolution}.  This proves the first-passage display and
fills the unexpanded calculation at source lines~1457--1463.

For the post-passage assertion, first sum the displayed first-passage bound
over its vertical coordinate.  The horizontal mass function is bounded by
\(C K_s(x-s/4)\), where
\[
 K_s(y)=(1+s)^{-1/2}
 \left(e^{-c y^2/(1+s)}+e^{-c|y|}\right).
\]
If \(\Sigma'\) is \(h\)-separated and \(r_0\le h/4\), the intervals of
radius \(r_0\) about its points are disjoint.  Since \(r_0\ge1\), direct
summation of the Gaussian and exponential pieces gives
\[
 \sum_{\operatorname{dist}(x,\Sigma')\le r_0}K_s(x-s/4)
 \le C\left(\frac {r_0}{\sqrt{1+s}}+\frac {r_0}h\right).
\]
Here \(h\le m^{0.4}=o(\sqrt{s})\), because
\(s>m/\log^2m\), so the right side is at most \(Cr_0/h\).  Conditional on the
sum of the \(p\) later horizontal increments, the target set is merely a
translate of \(\Sigma\) and has the same spacing.  Averaging the uniform
bound proves the assertion.  The condition \(h\le m^{0.4}\) is essential;
without it a singleton at the kernel peak would contradict a bare
\(O(r_0/h)\) claim.
\end{proof}

\begin{proposition}[repaired renewal monotonicity]
\label{prop:Tao-v7-renewal-monotonicity}
Retain \(n,\xi\) and the black--white partition from
Proposition~\ref{prop:Tao-v7-black-triangles}.  There is an absolute
\(\varepsilon_0>0\), fixed in the exit-rectangle calculation below, and we
assume \(0<\varepsilon\le\varepsilon_0\).  Extend \(1_W\) by zero off
\([n/2]\times\mathbb Z\).  For \(x\in\mathbb Z^2\), let
\[
 Q(x)=\mathbb E_x\prod_{q\ge0}e^{-\varepsilon^31_W(X_q)},
 \qquad X_q=x+\mathbf V_q,
\]
and define
\[
 Q_m=\sup_{\substack{j\in\Npos,\ j\ge\lfloor n/2\rfloor-m\\l\in\mathbb Z}}
 \max(\lfloor n/2\rfloor-j,1)^A Q(j,l).
\]
For every fixed sufficiently large \(A\), there is an explicit finite
threshold \(C_{A,\varepsilon}\) such that
\[
 Q_m\le Q_{m-1}\qquad
 (m\in\Npos,\ C_{A,\varepsilon}\le m\le\lfloor n/2\rfloor).
\]
The proof uses the weaker v7 Lemma~7.9 and does not assume
\eqref{eq:Tao-v5-rip}.
\end{proposition}

\begin{proof}
If \(m=\lfloor n/2\rfloor\), the restrictions defining \(Q_m\) and
\(Q_{m-1}\) both reduce to \(j\in\Npos\), so the asserted inequality is an
equality.  Hence suppose \(m\le\lfloor n/2\rfloor-1\).  It then suffices to
treat the genuine boundary row
\(j=\lfloor n/2\rfloor-m\in\Npos\).
If \((j,l)\) is white, the recursion
\[
 Q(j,l)=e^{-\varepsilon^3}\mathbb E Q((j,l)+\mathbf H)
\]
and
\[
 \max(m-r,1)^{-1}\le m^{-1}
 \exp\left(Cr\frac{\log m}{m}\right)
\]
show that the immediate white factor absorbs the horizontal exponential
moment for sufficiently large \(m\).  This is source Case~1.

Suppose next that \((j,l)\) lies at depth
\(0\le s\le m/\log^2m\) in a black triangle.  Stop at \(\kappa_s\) and
apply the recursion once more than the source does:
\[
 Q(j,l)=\mathbb E\left[
 e^{-\varepsilon^3\sum_{q=0}^{\kappa_s}1_W(X_q)}
 Q(X_{\kappa_s+1})\right].
\]
The definition of \(Q_{m-1}\) now applies at the terminal point.  With
\(\lambda=CA\log m/m\),
\[
 Q(j,l)\le m^{-A}Q_{m-1}
 \mathbb E\left[e^{-\varepsilon^31_W(X_{\kappa_s})}
 e^{\lambda(\mathbf V_{\kappa_s+1})_1}\right].
\]
Lemma~\ref{lem:Tao-v7-first-passage-repaired} gives
\(\mathbb E e^{\lambda(\mathbf V_{\kappa_s})_1}
=1+O_A(1/\log m)\).  We now prove, rather than assume, the uniform exit
probability.  Summing that lemma first over
\(l-s>M_v\), and then over
\(|j-s/4|>M_h\sqrt{1+s}\), gives
\[
 \mathbb P\left(
  1\le(\mathbf V_{\kappa_s})_2-s\le M_v,
  \ |(\mathbf V_{\kappa_s})_1-s/4|
       \le M_h\sqrt{1+s}
 \right)\ge\frac12
\]
once the two absolute integers \(M_v,M_h\) are chosen sufficiently large.
Write \(J=(\mathbf V_{\kappa_s})_1\) and
\(L=(\mathbf V_{\kappa_s})_2\).  In this rectangle
\(1\le L-s\le M_v\), so the exit point is above the top row by at most
\(M_v\).  Horizontally,
\[
 j+J-j_\Delta\ge \frac s4-M_h\sqrt{1+s}\ge-C_-(M_h),
\]
while the triangle inequality and \(J\le s/4+M_h\sqrt{1+s}\) give
\begin{align*}
 (j+J-j_\Delta)\log9
 &\le s_\Delta-s\log2+\frac{s\log9}{4}
       +M_h\log9\sqrt{1+s}\\
 &\le s_\Delta+C_+(M_h).
\end{align*}
The last constants are finite because
\(\log2-\frac14\log9>0\), so a negative linear term absorbs the square-root
term.  Thus the horizontal coordinate lies within an absolute distance of
the top-row segment
\([j_\Delta,j_\Delta+s_\Delta/\log9]\).  Combining the horizontal and
vertical bounds gives a Euclidean distance
\(C_{\rm exit}=C(M_v,M_h)\) from the triangle.  Choose the global
\(\varepsilon_0\le e^{-10}\) so that
\[
 \frac1{10}\log(1/\varepsilon_0)>C_{\rm exit}+1.
\]
The white collar then gives
\(\mathbb P(X_{\kappa_s}\in W)\ge c_*\) with \(c_*=1/2\), uniformly in
the triangle, its depth, \(n\), and \(\xi\).

Strong Markov factors the final holding increment.  Its exponential moment
is \(1+O_A(\log m/m)\).  Since the exponential weight accumulated before
that increment is at least one, the remaining expectation is at most
\[
 1+O_A(1/\log m)-c_*(1-e^{-\varepsilon^3}),
\]
before multiplication by the final \(1+o(1)\) factor, and is at most one for
\(m\ge C_{A,\varepsilon}\).  This repairs the unsupported source line~1556.

It remains to treat a starting point at depth
\(s>m/\log^2m\).  We first prove the two exact auxiliary bounds used in this
case.  Membership in the initial triangle and the real right-tip bound give
the exact upper range
\begin{equation}
\label{eq:Tao-deep-depth-upper}
 \frac m{\log^2m}<s\le\frac{\log9}{\log2}\,m.
\end{equation}
Put \(x=(j,l)\), retain the first-passage time \(\kappa_s\), and reset the
post-exit process by
\[
 \widetilde X_p=X_{\kappa_s+p}=x+\mathbf V_{\kappa_s+p},\qquad
 \widetilde{\mathbf V}_p
   =\mathbf V_{\kappa_s+p}-\mathbf V_{\kappa_s}
 \quad(p\in\Nzero).
\]
Conditionally on \(\mathcal F_{\kappa_s}\), the increments of
\(\widetilde{\mathbf V}\) are iid copies of \(\mathbf H\) and are
independent of \(\widetilde X_0\).

We first prove the v7 many-triangles estimate uniformly for a renewal process
\(Z_p=z+\mathbf V_p\) started at an arbitrary \(z\in\mathbb Z^2\).  Define
\(t_1\) as its first black time, allowing \(t_1=0\).  After entry into the
\(i\)-th triangle \(\Delta_i\), let \(\sigma_i>t_i\) be the first time whose
vertical coordinate exceeds \(l_{\Delta_i}\), and let \(t_{i+1}\ge\sigma_i\)
be the first subsequent black time.  Let \(r\in\Nzero\cup\{\infty\}\) be
the number of entries.  The exit-rectangle argument gives, on \(r\ge i\),
\[
 \mathbb P_z(Z_{\sigma_i}\in W\mid\mathcal F_{t_i})\ge c_*.
\]
Set \(\rho=1-(1-e^{-1})c_*\in(0,1)\), and decrease the fixed
\(\varepsilon_0\) once more so that \(e^{\varepsilon_0}\rho\le1\).  For
\(q\in\Nzero\), define the uniformly typed contraction
\[
 H_q(z)=\mathbb E_z\left[
  1_{\{r\ge q+1\}}\prod_{i=1}^{q}e^{-1_W(Z_{\sigma_i})}
 \right].
\]
Since \(H_0(z)\le1\), strong Markov at the first exit gives inductively
\begin{align*}
 H_q(z)
 &=\mathbb E_z\left[
  1_{\{r\ge1\}}e^{-1_W(Z_{\sigma_1})}
  H_{q-1}(Z_{\sigma_1})\right]\\
 &\le\rho^{q-1}\mathbb E_z\left[
  1_{\{r\ge1\}}e^{-1_W(Z_{\sigma_1})}\right]
 \le\rho^q.
\end{align*}
On \(r\ge R\), the distinct exits
\(\sigma_1,\ldots,\sigma_{R-1}\) occur no later than \(t_R\); hence the
complete white count dominates their indicators and
\begin{equation}
\label{eq:Tao-v7-many-triangles-proved}
 \mathbb E_z1_{\{r\ge R\}}
 e^{-\sum_{p=1}^{t_R}1_W(Z_p)+\varepsilon R}
 \le e^{\varepsilon R}\rho^{R-1}\le e^\varepsilon.
\end{equation}
This proves the weaker v7 Lemma~7.9 with a defined time-zero atom and the
required uniform induction.

We next reconstruct the parameterised large-triangle estimate.  For
\(p\in\Nzero\) and \(s'>0\), define the event
\[
 E_{p,s'}=\{\widetilde X_p\in\Delta'
   \text{ for some }\Delta'\in\mathcal T
   \text{ with }s_{\Delta'}\ge s'\}.
\]
Let \(D\ge D_0\), where \(D_0\) is a sufficiently large absolute constant.
The first-passage overshoot bound and the exponential tail of the next
\(p\) increments give
\[
 \mathbb P\bigl((\widetilde X_p)_2
     \ge l_\Delta+2D(1+p)\bigr)
 \le C e^{-cD(1+p)}.
\]
The horizontal first-passage kernel and the same increment tail, using
\(p\le m^{0.1}\), \(s>m/\log^2m\), and \(s'\le m^{0.4}\), give
\[
 \mathbb P\left(
 \left|(\mathbf V_{\kappa_s+p})_1-\frac s4\right|>2s^{0.6}
 \right)
 \le C\frac{D(1+p)}{s'}.
\]
Outside these two exceptional events, membership of \(\widetilde X_p\) in
an eligible \(\Delta'\) forces its lower tip to satisfy
\begin{equation}
\label{eq:Tao-lower-tip-localisation}
 l_\Delta-10<l_{\Delta'}-\frac{s_{\Delta'}}{\log2}
 \le l_\Delta+2D(1+p).
\end{equation}
For if the left inequality failed, moving horizontally by
\(O(D(1+p))\) on the row \(l_\Delta\) would produce a lattice point of
\(\Delta'\); the horizontal bound and
\(\frac14\log9<\log2\) put the same point in the initial \(\Delta\),
contradicting disjointness.  The right inequality follows directly from
the triangle inequality for \(\widetilde X_p\in\Delta'\).  Substituting
\eqref{eq:Tao-lower-tip-localisation} back into that triangle inequality
gives
\[
 |(\widetilde X_p)_1-j_{\Delta'}|\le C D(1+p).
\]
If \(\Delta',\Delta''\) are two eligible triangles, intersecting each with
the row \(l_\Delta+\lfloor s'/2\rfloor\) produces disjoint horizontal
integer intervals of length at least
\(c s'-CD(1+p)\).  Thus, provided
\(s'\ge C_1D(1+p)\), their top-left coordinates are at least \(c_1s'\)
apart.  Decrease the separation constant once, if necessary, so that
\(0<c_1\le1\).  Apply the post-passage part of
Lemma~\ref{lem:Tao-v7-first-passage-repaired} with neighbourhood radius
\(r_0=CD(1+p)\) and spacing \(h=c_1s'\).  Enlarging \(C_1\) enforces
\(4r_0\le h\), and \(s'\le m^{0.4}\) enforces the smoothing-scale bound.
Consequently there are absolute \(C_0,c_0,C_1>0\) such that
\begin{equation}
\label{eq:Tao-large-triangle-parametric}
 \mathbb P(E_{p,s'})\le C_0\left(
 \frac{D(1+p)}{s'}+e^{-c_0D(1+p)}\right)
\end{equation}
whenever
\[
 0\le p\le m^{0.1},\qquad
 C_1D(1+p)\le s'\le m^{0.4}.
\]
This exposes the lower-tip and separated-top steps preceding the
post-\(p\) convolution at source lines~1805--1812.

Now put \(\delta_A=10^{-A-2}\), and choose
\[
 D_A\ge\max\left(D_0,
 \frac1{c_0}\log\frac{6C_0}{\delta_A(1-e^{-c_0})}\right),
 \qquad
 K_A\ge\max\left(C_1D_A,
 \frac{6C_0D_A\zeta(2)}{\delta_A}\right).
\]
Set \(H=10A/\varepsilon^3\) and
\[
 R=\left\lceil
 \frac{H+\log3+(A+2)\log10+\varepsilon+1}{\varepsilon}
 \right\rceil.
\]
Let \(h=\lceil H\rceil\), put \(u_1=h\), and define, for
\(1\le i<R\),
\[
 u_{i+1}=u_i+\left\lceil10K_A(1+u_i)^3\right\rceil+h+1,
 \qquad P=u_R+2.
\]
All of \(D_A,K_A,H,R,h,u_i,P\) depend only on \(A,\varepsilon\).  Enlarge
\(C_{A,\varepsilon}\) so that
\begin{equation}
\label{eq:Tao-case3-applicability}
 P\le m^{0.1},\qquad K_A(1+P)^3\le m^{0.4}.
\end{equation}

Only now define the finite exceptional event
\[
 E_*^{(P)}=\bigcup_{p=0}^{P-1}E_{p,K_A(1+p)^3}.
\]
Every term lies in the domain of
\eqref{eq:Tao-large-triangle-parametric}, and the union bound gives
\begin{align}
 \mathbb P(E_*^{(P)})
 &\le \frac{C_0D_A}{K_A}
       \sum_{p=0}^{P-1}\frac1{(1+p)^2}
   +C_0\sum_{p=0}^{P-1}e^{-c_0D_A(1+p)}\notag\\
 &\le\delta_A/6+\delta_A/6=\delta_A/3.
 \label{eq:Tao-E-star-budget}
\end{align}
This replaces the source's insufficient
\(O(A^2 4^{-A})\): its ratio to \(10^{-A-2}\) is
\(100A^2(10/4)^A\), which is unbounded.

Condition on \(\widetilde X_0\) and let \(r,t_i,\sigma_i\) be the preceding
triangle statistics for the shifted process \((\widetilde X_p)_{p\ge0}\).
Put
\[
 S_R=\sum_{p=1}^{t_R}1_W(\widetilde X_p),\qquad
 F_*=\left\{r\ge R,
 e^{-S_R+\varepsilon R}>3\delta_A^{-1}e^\varepsilon\right\}.
\]
The uniform estimate \eqref{eq:Tao-v7-many-triangles-proved}, integrated in
\(\widetilde X_0\), and Markov's inequality give
\(\mathbb P(F_*)\le\delta_A/3\).  Outside \(F_*\), the event \(r\ge R\)
forces
\[
 S_R\ge\varepsilon R-\log3-\log(\delta_A^{-1})-\varepsilon>H.
\]

Define the high-horizontal event
\[
 G=\{(\mathbf V_{\kappa_s+P})_1\ge0.9m\}.
\]
By \eqref{eq:Tao-deep-depth-upper} and the first-passage kernel,
\[
 \mathbb P((\mathbf V_{\kappa_s})_1\ge0.8m)\le Ce^{-cm},
\]
because \(0.8>\frac14\frac{\log9}{\log2}\).  Strong Markov and the
exponential tail of the fixed sum
\((\widetilde{\mathbf V}_P)_1\) give
\[
 \mathbb P(G)\le C_Pe^{-cm}.
\]
Enlarge \(C_{A,\varepsilon}\) again so that
\begin{equation}
\label{eq:Tao-G-budget}
 \mathbb P(G)\le\delta_A/3,
 \qquad m^A\mathbb P(G)\le\frac12.
\end{equation}
Outside \(G\), positivity of the horizontal increments and
\(j=\lfloor n/2\rfloor-m\) keep every
\(\widetilde X_p\), \(0\le p<P\), inside the strip where
\(B\uplus W\) is the full partition.

Define the shifted white count
\[
 N_W=\sum_{p=0}^{P-1}1_W(\widetilde X_p).
\]
Suppose \(N_W\le H\) and none of \(E_*^{(P)},F_*,G\) occurs.  Among the
\(h+1\) indices \(0,\ldots,h\) at least one is black.  Therefore the actual
first entry time defined above satisfies \(t_1\le h=u_1\).  Induct on the
recursively defined entry times, not on an arbitrary selection of black
points.  If \(t_i\le u_i\), then, outside \(E_*^{(P)}\), its triangle
\(\Delta_i\) satisfies
\[
 s_{\Delta_i}<K_A(1+t_i)^3\le K_A(1+u_i)^3.
\]
Every holding increment has vertical component at least three, so by time
\[
 b_i=t_i+\left\lceil10K_A(1+u_i)^3\right\rceil+1
\]
the process is strictly above \(l_{\Delta_i}\), and in particular
\(b_i\ge\sigma_i\).  The block \(b_i,\ldots,b_i+h\) lies in
\(\{0,\ldots,P-1\}\) and contains a black point; otherwise that block alone
would contribute \(h+1>H\) white points.  By the defining minimality of the
next entry time,
\[
 t_{i+1}\le b_i+h
 \le u_i+\left\lceil10K_A(1+u_i)^3\right\rceil+h+1
 =u_{i+1}.
\]
Thus \(t_{i+1}\) exists and is the entry into the next distinct triangle.
Iteration gives \(R\) entries and \(t_R\le u_R\le P-2\).  Hence
\(S_R\le N_W\le H\), contradicting the preceding consequence outside
\(F_*\).  Therefore
\begin{equation}
\label{eq:Tao-low-white-budget}
 \mathbb P(N_W\le H)
 \le\mathbb P(E_*^{(P)})+\mathbb P(F_*)+\mathbb P(G)
 \le\delta_A.
\end{equation}

It remains to connect this probability estimate to the recursion.  Dropping
the factors before \(\kappa_s\) and then iterating exactly \(P\) times gives
\begin{align*}
 Q(j,l)
 &\le\mathbb E\left[e^{-\varepsilon^3N_W}Q(\widetilde X_P)\right]\\
 &\le m^{-A}Q_{m-1}\,
 \mathbb E\left[e^{-\varepsilon^3N_W}
 \max\left(1-\frac{(\mathbf V_{\kappa_s+P})_1}{m},
                 \frac1m\right)^{-A}\right].
\end{align*}
The second line uses the fact that the terminal horizontal coordinate has
advanced by at least one, so the defining range of \(Q_{m-1}\) applies.
On \(G\) the weight is at most \(m^A\); on \(G^c\) it is at most \(10^A\).
Moreover,
\[
 \mathbb E e^{-\varepsilon^3N_W}
 \le\mathbb P(N_W\le H)+e^{-\varepsilon^3H}
 \le\delta_A+e^{-10A}\le10^{-A-1}
\]
for sufficiently large \(A\).  By \eqref{eq:Tao-G-budget}, the complete
weighted expectation is at most
\[
 m^A\mathbb P(G)+10^A(\delta_A+e^{-10A})\le1
\]
after the stated enlargement of \(C_{A,\varepsilon}\).  This proves
\(Q_m\le Q_{m-1}\) in the final case.
\end{proof}

\begin{corollary}[closure of the v7 fine-scale branch]
\label{cor:Tao-v7-fine-scale-branch}
With the repairs proved above,
\[
 Q(j,l)\ll_A\max(\lfloor n/2\rfloor-j,1)^{-A}
 \quad(j\in\Npos,\ l\in\mathbb Z),
 \qquad \mathbb E Q(\mathbf H)\ll_A n^{-A}.
\]
Consequently the v7 implications
\[
 \text{Proposition 7.8}\Rightarrow
 \text{Proposition 7.3}\Rightarrow
 \text{Proposition 7.1}\Rightarrow
 \text{Proposition 1.17}\Rightarrow
 \text{Proposition 1.14}
\]
hold at their printed uniform scope.  The exact downstream dependency in the
source is
\[
 \text{Proposition 1.14}+\text{Proposition 1.9}
 \Rightarrow\text{Proposition 1.11}
 \Rightarrow\text{Theorem 1.6}
 \Rightarrow\text{Theorem 1.3}.
\]
\end{corollary}

\begin{proof}
Monotonicity and the finite base range bound all \(Q_m\) independently of
\(m\).  The first coordinate of \(\mathbf H\) has exponential tail, so
splitting at \(\mathbf H_1=n/4\) gives
\(\mathbb E Q(\mathbf H)=O_A(n^{-A})\).  This is Proposition~7.3.
The white-point pairing in
Proposition~\ref{prop:Tao-v7-pairing-holding} gives Proposition~7.1, whose
reversed-variable form is Proposition~1.17.  Finally
Proposition~\ref{prop:Tao-v7-Fourier-to-mixing} gives Proposition~1.14.
The second display records every consequence edge leaving this branch.
Proposition~\ref{prop:Tao-v7-Prop19-valuation-law} now supplies the separately
typed Proposition~1.9 input.  The exact step combining it with
Proposition~1.14 to prove Proposition~1.11, and the two later theorem edges,
remain behind their named proof-audit gate.  This is a dependency record, not
a handwaving theorem claim.
\end{proof}

\begin{nonclaim}
Corollary~\ref{cor:Tao-v7-fine-scale-branch} certifies this versioned proof
branch; it does not replace the durable whole-paper release audit.  By itself
it does not certify Proposition~1.9, Proposition~1.11, Theorem~1.6, or
Theorem~1.3; Proposition~\ref{prop:Tao-v7-Prop19-valuation-law} supplies the
first of those results by an independent argument.  This branch does
not prove the stronger v5/journal Lemma~7.9, use Baker's theorem, establish
exponential Fourier decay, prove natural-density mixing of the full affine
map, give a fixed absolute orbit bound, or prove convergence to one.
\end{nonclaim}
